\section{Asynchronous COM-Frame Streamfunction Fit with Fast COM/Background Basis and Slow Mesoscale Basis}

We consider an asynchronous Lagrangian dataset in which the drifters are observed over a
common time window $[t_0,t_f]$, but not necessarily at simultaneous times. The total velocity
of each drifter is modeled as the sum of a mesoscale component, a background component, and a
submesoscale component. The mesoscale component is represented by a tensor-product
streamfunction with a slow temporal spline basis, while the center-of-mass trajectory and the
background flow are represented using a faster temporal spline basis constructed directly from
the pooled observation times.

\subsection{Asynchronous observations}

Assume that drifter $k$ is observed at times
\begin{equation}
t_{k,1},\dots,t_{k,N_k}, \qquad k=1,\dots,K,
\end{equation}
all lying in a common interval $[t_0,t_f]$, but not necessarily with
$t_{k,n}=t_{k',n'}$ for different drifters.

Let the observed positions be
\begin{equation}
x_{k,n} \equiv x_k(t_{k,n}),
\qquad
y_{k,n} \equiv y_k(t_{k,n}),
\qquad
n=1,\dots,N_k,
\end{equation}
and let the corresponding observed velocities be
\begin{equation}
\dot{x}_{k,n} \equiv \dot{x}_k(t_{k,n}),
\qquad
\dot{y}_{k,n} \equiv \dot{y}_k(t_{k,n}),
\end{equation}
obtained either directly or from a separate trajectory-smoothing step.

\subsection{Total velocity model}

We write the full model in the fixed frame as
\begin{equation}
\dot{x}_k(t)=u_k^{\mathrm{meso}}(t)+u^{\mathrm{bg}}(t)+u_k^{\mathrm{sm}}(t),
\qquad
\dot{y}_k(t)=v_k^{\mathrm{meso}}(t)+v^{\mathrm{bg}}(t)+v_k^{\mathrm{sm}}(t),
\label{eq:total_model_full}
\end{equation}
where
\begin{itemize}
    \item $u_k^{\mathrm{meso}},v_k^{\mathrm{meso}}$ are the mesoscale velocities induced by a
    streamfunction,
    \item $u^{\mathrm{bg}},v^{\mathrm{bg}}$ are the background velocities, common to all drifters
    and depending only on time,
    \item $u_k^{\mathrm{sm}},v_k^{\mathrm{sm}}$ are the remaining submesoscale velocities.
\end{itemize}

During the core mesoscale fit, however, we do \emph{not} attempt to separate the background
and submesoscale parts. Instead, we combine them into a single residual:
\begin{equation}
r_k^x(t)=u^{\mathrm{bg}}(t)+u_k^{\mathrm{sm}}(t),
\qquad
r_k^y(t)=v^{\mathrm{bg}}(t)+v_k^{\mathrm{sm}}(t).
\label{eq:combined_residual}
\end{equation}
Thus the total model used during the mesoscale computation is
\begin{equation}
\dot{x}_k(t)=u_k^{\mathrm{meso}}(t)+r_k^x(t),
\qquad
\dot{y}_k(t)=v_k^{\mathrm{meso}}(t)+r_k^y(t).
\label{eq:core_total_model}
\end{equation}

The background and submesoscale parts are separated only after the mesoscale fit has been
completed.

\subsection{Fast temporal basis}

Let the full collection of observation times, pooled across all drifters and sorted in
nondecreasing order, be
\begin{equation}
\tau_1 \le \tau_2 \le \cdots \le \tau_{N_{\mathrm{tot}}},
\qquad
N_{\mathrm{tot}}=\sum_{k=1}^{K}N_k.
\end{equation}
Let
\begin{equation}
D(\tau_r)
\end{equation}
denote the number of deployed drifters at time $\tau_r$, that is, the number of drifters whose
lifetimes include $\tau_r$.

We define a subsequence of representative times
\begin{equation}
\sigma_1,\sigma_2,\dots,\sigma_Q
\end{equation}
recursively from the pooled sorted times by
\begin{equation}
q_1=1,
\qquad
q_{m+1}=q_m + D(\tau_{q_m}),
\label{eq:fast_stride_rule}
\end{equation}
until $q_m>N_{\mathrm{tot}}$, and then set
\begin{equation}
\sigma_m=\tau_{q_m}, \qquad m=1,\dots,Q.
\label{eq:fast_knot_times}
\end{equation}

These representative times are used as knot locations for a fast temporal spline basis
\begin{equation}
\{B_q(t)\}_{q=1}^{Q_b}.
\label{eq:fast_basis_def}
\end{equation}
This fast basis will be used
\begin{enumerate}
    \item to represent the center-of-mass trajectory, and
    \item to represent the background flow.
\end{enumerate}

\subsection{Slow temporal basis for the mesoscale streamfunction}

Let
\begin{equation}
\{T_\ell(t)\}_{\ell=1}^{L}
\label{eq:slow_basis_def}
\end{equation}
be a slow temporal spline basis on $[t_0,t_f]$. This basis is used only for the mesoscale
streamfunction.

\subsection{Spline representation of the center of mass in the fast basis}

We represent the center-of-mass coordinates as
\begin{equation}
m_x(t)=\sum_{q=1}^{Q_b} b^x_q B_q(t),
\qquad
m_y(t)=\sum_{q=1}^{Q_b} b^y_q B_q(t),
\label{eq:com_fast_basis}
\end{equation}
with unknown coefficients $b^x_q$ and $b^y_q$.

At each drifter observation time $t_{k,n}$, define
\begin{equation}
m_{x;k,n}=m_x(t_{k,n})=\sum_{q=1}^{Q_b} b^x_q B_q(t_{k,n}),
\qquad
m_{y;k,n}=m_y(t_{k,n})=\sum_{q=1}^{Q_b} b^y_q B_q(t_{k,n}).
\end{equation}

The corresponding COM-frame coordinates are
\begin{equation}
\tilde{x}_{k,n}=x_{k,n}-m_{x;k,n},
\qquad
\tilde{y}_{k,n}=y_{k,n}-m_{y;k,n}.
\label{eq:relative_coords_fastcom}
\end{equation}

Differentiating \eqref{eq:com_fast_basis} gives the COM velocity
\begin{equation}
\dot{m}_x(t)=\sum_{q=1}^{Q_b} b^x_q \dot{B}_q(t),
\qquad
\dot{m}_y(t)=\sum_{q=1}^{Q_b} b^y_q \dot{B}_q(t).
\label{eq:com_velocity_fastcom}
\end{equation}

\subsection{Least-squares fit for the COM coefficients}

We estimate the COM coefficients by fitting a common smooth trajectory to all drifter
positions:
\begin{align}
\min_{\{b^x_q\}}
\sum_{k=1}^{K}\sum_{n=1}^{N_k}
\left(
x_{k,n}-\sum_{q=1}^{Q_b} b^x_q B_q(t_{k,n})
\right)^2,
\\
\min_{\{b^y_q\}}
\sum_{k=1}^{K}\sum_{n=1}^{N_k}
\left(
y_{k,n}-\sum_{q=1}^{Q_b} b^y_q B_q(t_{k,n})
\right)^2.
\end{align}

Stack all observations using a single index $r=1,\dots,N_{\mathrm{tot}}$, with associated
pairs $(k(r),n(r))$, and define
\begin{equation}
\mathbf{x}=
\begin{bmatrix}
x_{k(r),n(r)}
\end{bmatrix}_{r=1}^{N_{\mathrm{tot}}},
\qquad
\mathbf{y}=
\begin{bmatrix}
y_{k(r),n(r)}
\end{bmatrix}_{r=1}^{N_{\mathrm{tot}}},
\end{equation}
together with the fast temporal design matrix
\begin{equation}
\mathcal{B}_{rq}=B_q\!\left(t_{k(r),n(r)}\right).
\label{eq:B_design_matrix}
\end{equation}
Then
\begin{equation}
\hat{\mathbf{b}}^x=(\mathcal{B}^\top\mathcal{B})^{-1}\mathcal{B}^\top\mathbf{x},
\qquad
\hat{\mathbf{b}}^y=(\mathcal{B}^\top\mathcal{B})^{-1}\mathcal{B}^\top\mathbf{y},
\label{eq:com_fast_coeff_est}
\end{equation}
or, with regularization,
\begin{equation}
\hat{\mathbf{b}}^x=(\mathcal{B}^\top\mathcal{B}+\lambda_c R_c)^{-1}\mathcal{B}^\top\mathbf{x},
\qquad
\hat{\mathbf{b}}^y=(\mathcal{B}^\top\mathcal{B}+\lambda_c R_c)^{-1}\mathcal{B}^\top\mathbf{y}.
\end{equation}

After estimating $\hat{\mathbf{b}}^x,\hat{\mathbf{b}}^y$, we evaluate
$\hat{m}_x(t_{k,n})$, $\hat{m}_y(t_{k,n})$ and form the COM-frame coordinates
$\tilde{x}_{k,n},\tilde{y}_{k,n}$ from \eqref{eq:relative_coords_fastcom}.

\subsection{Tensor-product mesoscale streamfunction in the COM frame}

Let $X_i(\tilde{x})$, $Y_j(\tilde{y})$, and $T_\ell(t)$ denote basis functions in
$\tilde{x}$, $\tilde{y}$, and $t$, respectively. We represent the mesoscale streamfunction as
\begin{equation}
\psi(\tilde{x},\tilde{y},t)
=
\sum_{i=1}^{n_x}\sum_{j=1}^{n_y}\sum_{\ell=1}^{L}
a_{ij\ell}\,
X_i(\tilde{x})\,Y_j(\tilde{y})\,T_\ell(t).
\label{eq:psi_slow_meso}
\end{equation}

The mesoscale velocity is
\begin{equation}
u^{\mathrm{meso}}=-\partial_{\tilde{y}}\psi,
\qquad
v^{\mathrm{meso}}=\partial_{\tilde{x}}\psi,
\end{equation}
hence
\begin{align}
u^{\mathrm{meso}}(\tilde{x},\tilde{y},t)
&=
-\sum_{i=1}^{n_x}\sum_{j=1}^{n_y}\sum_{\ell=1}^{L}
a_{ij\ell}\,
X_i(\tilde{x})\,Y_j'(\tilde{y})\,T_\ell(t),
\label{eq:u_meso_slow}
\\
v^{\mathrm{meso}}(\tilde{x},\tilde{y},t)
&=
\sum_{i=1}^{n_x}\sum_{j=1}^{n_y}\sum_{\ell=1}^{L}
a_{ij\ell}\,
X_i'(\tilde{x})\,Y_j(\tilde{y})\,T_\ell(t).
\label{eq:v_meso_slow}
\end{align}

Evaluated at the drifter observations,
\begin{align}
u^{\mathrm{meso}}_{k,n}
&=
-\sum_{i,j,\ell}
a_{ij\ell}\,
X_i(\tilde{x}_{k,n})\,Y_j'(\tilde{y}_{k,n})\,T_\ell(t_{k,n}),
\label{eq:u_meso_kn}
\\
v^{\mathrm{meso}}_{k,n}
&=
\sum_{i,j,\ell}
a_{ij\ell}\,
X_i'(\tilde{x}_{k,n})\,Y_j(\tilde{y}_{k,n})\,T_\ell(t_{k,n}).
\label{eq:v_meso_kn}
\end{align}

\subsection{Core computation: combine background and submesoscale into a single residual}

Differentiating \eqref{eq:relative_coords_fastcom} gives
\begin{equation}
\dot{\tilde{x}}_k(t)=\dot{x}_k(t)-\dot{m}_x(t),
\qquad
\dot{\tilde{y}}_k(t)=\dot{y}_k(t)-\dot{m}_y(t).
\end{equation}
During the core mesoscale fit, we write
\begin{align}
\dot{\tilde{x}}_{k,n}
&=
-\sum_{i,j,\ell}
a_{ij\ell}\,
X_i(\tilde{x}_{k,n})\,Y_j'(\tilde{y}_{k,n})\,T_\ell(t_{k,n})
+\eta^x_{k,n},
\label{eq:core_residual_x}
\\
\dot{\tilde{y}}_{k,n}
&=
\sum_{i,j,\ell}
a_{ij\ell}\,
X_i'(\tilde{x}_{k,n})\,Y_j(\tilde{y}_{k,n})\,T_\ell(t_{k,n})
+\eta^y_{k,n},
\label{eq:core_residual_y}
\end{align}
where the combined residual satisfies
\begin{equation}
\eta^x_{k,n}
=
u^{\mathrm{bg}}(t_{k,n})+u^{\mathrm{sm}}_{k,n},
\qquad
\eta^y_{k,n}
=
v^{\mathrm{bg}}(t_{k,n})+v^{\mathrm{sm}}_{k,n}.
\label{eq:eta_decomp}
\end{equation}

Thus the mesoscale fit is performed without yet distinguishing background and submesoscale
motions.

\subsection{Linear least-squares problem for the mesoscale coefficients}

Let
\begin{equation}
P=n_x n_y L
\end{equation}
be the total number of mesoscale streamfunction coefficients, and flatten $(i,j,\ell)$ to a
single column index $\alpha=1,\dots,P$. Define
\begin{equation}
\mathbf{a}=
\begin{bmatrix}
a_{111}\\
a_{211}\\
\vdots\\
a_{n_x n_y L}
\end{bmatrix}.
\end{equation}

Stack all observed COM-frame velocities:
\begin{equation}
\mathbf{u}=
\begin{bmatrix}
\dot{\tilde{x}}_{k,n}\\
\dot{\tilde{y}}_{k,n}
\end{bmatrix}_{(k,n)},
\qquad
\boldsymbol{\eta}=
\begin{bmatrix}
\eta^x_{k,n}\\
\eta^y_{k,n}
\end{bmatrix}_{(k,n)}.
\end{equation}

For coefficient $a_{ij\ell}$, define the corresponding design-matrix column
\begin{equation}
\mathbf{X}_{:,\,ij\ell}
=
\begin{bmatrix}
-\,X_i(\tilde{x}_{k,n})\,Y_j'(\tilde{y}_{k,n})\,T_\ell(t_{k,n})
\\[0.25em]
\phantom{-}\,X_i'(\tilde{x}_{k,n})\,Y_j(\tilde{y}_{k,n})\,T_\ell(t_{k,n})
\end{bmatrix}_{(k,n)}.
\label{eq:X_design_meso}
\end{equation}
Then
\begin{equation}
\mathbf{u}=\mathbf{X}\mathbf{a}+\boldsymbol{\eta}.
\label{eq:meso_regression}
\end{equation}

The least-squares estimator is
\begin{equation}
\hat{\mathbf{a}}
=
(\mathbf{X}^\top\mathbf{X})^{-1}\mathbf{X}^\top\mathbf{u},
\end{equation}
or, with regularization,
\begin{equation}
\hat{\mathbf{a}}
=
(\mathbf{X}^\top\mathbf{X}+\lambda_\psi R_\psi)^{-1}\mathbf{X}^\top\mathbf{u}.
\label{eq:meso_regression_regularized}
\end{equation}

\subsection{Recover the combined residual after the mesoscale fit}

Once $\hat{\mathbf{a}}$ is known, the fitted mesoscale velocities are
\begin{align}
\hat{u}^{\mathrm{meso}}_{k,n}
&=
-\sum_{i,j,\ell}
\hat{a}_{ij\ell}\,
X_i(\tilde{x}_{k,n})\,Y_j'(\tilde{y}_{k,n})\,T_\ell(t_{k,n}),
\\
\hat{v}^{\mathrm{meso}}_{k,n}
&=
\sum_{i,j,\ell}
\hat{a}_{ij\ell}\,
X_i'(\tilde{x}_{k,n})\,Y_j(\tilde{y}_{k,n})\,T_\ell(t_{k,n}).
\end{align}

The combined residual in the fixed frame is then
\begin{equation}
\rho^x_{k,n}=\dot{x}_{k,n}-\hat{u}^{\mathrm{meso}}_{k,n},
\qquad
\rho^y_{k,n}=\dot{y}_{k,n}-\hat{v}^{\mathrm{meso}}_{k,n}.
\label{eq:rho_fixed_residual}
\end{equation}
By construction, this residual contains both the background and the submesoscale parts:
\begin{equation}
\rho^x_{k,n}=u^{\mathrm{bg}}(t_{k,n})+u^{\mathrm{sm}}_{k,n},
\qquad
\rho^y_{k,n}=v^{\mathrm{bg}}(t_{k,n})+v^{\mathrm{sm}}_{k,n}.
\label{eq:rho_decomposition}
\end{equation}

\subsection{Separate the background from the combined residual}

We now define the background as the projection of the residual mean velocity onto the fast
basis $B_q(t)$. Specifically,
\begin{equation}
u^{\mathrm{bg}}(t)=\sum_{q=1}^{Q_b} c^x_q B_q(t),
\qquad
v^{\mathrm{bg}}(t)=\sum_{q=1}^{Q_b} c^y_q B_q(t),
\label{eq:bg_fast_rep}
\end{equation}
with coefficients determined by
\begin{align}
\hat{\mathbf{c}}^x
&=
\arg\min_{\mathbf{c}^x}
\sum_{k=1}^{K}\sum_{n=1}^{N_k}
\left(
\rho^x_{k,n}-\sum_{q=1}^{Q_b} c^x_q B_q(t_{k,n})
\right)^2,
\label{eq:bg_proj_x}
\\
\hat{\mathbf{c}}^y
&=
\arg\min_{\mathbf{c}^y}
\sum_{k=1}^{K}\sum_{n=1}^{N_k}
\left(
\rho^y_{k,n}-\sum_{q=1}^{Q_b} c^y_q B_q(t_{k,n})
\right)^2.
\label{eq:bg_proj_y}
\end{align}

Equivalently, using the fast design matrix \eqref{eq:B_design_matrix} and the stacked
residuals
\begin{equation}
\boldsymbol{\rho}^x=
\begin{bmatrix}
\rho^x_{k(r),n(r)}
\end{bmatrix}_{r=1}^{N_{\mathrm{tot}}},
\qquad
\boldsymbol{\rho}^y=
\begin{bmatrix}
\rho^y_{k(r),n(r)}
\end{bmatrix}_{r=1}^{N_{\mathrm{tot}}},
\end{equation}
we have
\begin{equation}
\hat{\mathbf{c}}^x=(\mathcal{B}^\top\mathcal{B})^{-1}\mathcal{B}^\top\boldsymbol{\rho}^x,
\qquad
\hat{\mathbf{c}}^y=(\mathcal{B}^\top\mathcal{B})^{-1}\mathcal{B}^\top\boldsymbol{\rho}^y,
\label{eq:bg_proj_closed_form}
\end{equation}
or, with regularization,
\begin{equation}
\hat{\mathbf{c}}^x=(\mathcal{B}^\top\mathcal{B}+\lambda_b R_b)^{-1}\mathcal{B}^\top\boldsymbol{\rho}^x,
\qquad
\hat{\mathbf{c}}^y=(\mathcal{B}^\top\mathcal{B}+\lambda_b R_b)^{-1}\mathcal{B}^\top\boldsymbol{\rho}^y.
\end{equation}

\subsection{Define the submesoscale component}

The submesoscale signal is defined as the residual that remains after subtracting the
background from the combined residual:
\begin{equation}
\hat{u}^{\mathrm{sm}}_{k,n}
=
\rho^x_{k,n}-\hat{u}^{\mathrm{bg}}(t_{k,n}),
\qquad
\hat{v}^{\mathrm{sm}}_{k,n}
=
\rho^y_{k,n}-\hat{v}^{\mathrm{bg}}(t_{k,n}),
\label{eq:submeso_final}
\end{equation}
or, equivalently,
\begin{equation}
\hat{u}^{\mathrm{sm}}_{k,n}
=
\dot{x}_{k,n}
-\hat{u}^{\mathrm{meso}}_{k,n}
-\hat{u}^{\mathrm{bg}}(t_{k,n}),
\qquad
\hat{v}^{\mathrm{sm}}_{k,n}
=
\dot{y}_{k,n}
-\hat{v}^{\mathrm{meso}}_{k,n}
-\hat{v}^{\mathrm{bg}}(t_{k,n}).
\end{equation}

\subsection{Consistency with the synchronous fixed-$K$ case}

Suppose all drifters are observed synchronously at times
\begin{equation}
t_1,\dots,t_N,
\end{equation}
and the number of deployed drifters is constant and equal to $K$. Then the pooled sorted
time list consists of each common observation time repeated $K$ times, so the stride rule
\eqref{eq:fast_stride_rule} selects one representative time per synchronous sample. Hence the
fast-basis knots coincide with the synchronous sampling times, up to spline-order details.

In this case, projecting the pooled residuals onto the common time-only basis $B_q(t)$ is
equivalent to projecting the synchronous drifter-mean residual velocity onto that basis.
Thus the present construction reduces to the Oscroft-style background definition when the data
are synchronous and the number of deployed drifters does not change.

\subsection{Compact three-stage algorithm}

The computation proceeds as follows.

\paragraph{Step 1.}
Pool and sort all observation times, construct the representative times
$\sigma_m$ from \eqref{eq:fast_stride_rule}--\eqref{eq:fast_knot_times}, and build the fast
basis $B_q(t)$.

\paragraph{Step 2.}
Fit the COM coefficients $b_q^x,b_q^y$ in the fast basis using
\eqref{eq:com_fast_coeff_est}, and compute the COM-frame coordinates
$\tilde{x}_{k,n},\tilde{y}_{k,n}$.

\paragraph{Step 3.}
Choose the slow mesoscale basis $T_\ell(t)$ and the spatial bases $X_i,Y_j$, assemble the
design matrix \eqref{eq:X_design_meso}, and solve
\eqref{eq:meso_regression} or \eqref{eq:meso_regression_regularized} for the mesoscale
coefficients $a_{ij\ell}$.

\paragraph{Step 4.}
Compute the fixed-frame combined residual \eqref{eq:rho_fixed_residual}.

\paragraph{Step 5.}
Project that residual onto the fast basis using \eqref{eq:bg_proj_x}--\eqref{eq:bg_proj_closed_form}
to obtain the background flow.

\paragraph{Step 6.}
Define the remaining residual by \eqref{eq:submeso_final} as the submesoscale component.

\subsection{Remarks}

\begin{enumerate}
    \item The fast basis $B_q(t)$ is used for both the COM fit and the background fit.
    \item The slow basis $T_\ell(t)$ is used only for the mesoscale streamfunction.
    \item The mesoscale fit remains linear because the background and submesoscale components
    are grouped together into a single residual during the core computation.
    \item The separation between background and submesoscale is performed only after the
    mesoscale coefficients have been estimated.
    \item In the synchronous fixed-$K$ limit, the present construction reduces to the
    corresponding Oscroft-style mean-residual definition.
\end{enumerate}

\subsection{Addendum: numerical implementation details for GriddedStreamfunction}

This addendum records the numerical choices used by the current MATLAB implementation
`GriddedStreamfunction.fitFromTrajectorySplines(...)`. These choices are sufficient to
reproduce the shipped algorithm exactly, but some of them are implementation defaults rather
than mathematically required parts of the formulation above.

\subsubsection{Implemented least-squares path}

The current implementation follows the three-stage least-squares structure above, but uses the
unregularized branch throughout. In practice,
\begin{enumerate}
    \item the COM coefficients are obtained from the pooled fast-basis system
    \begin{equation}
    \hat{\mathbf{b}}^x = \mathcal{B} \backslash \mathbf{x},
    \qquad
    \hat{\mathbf{b}}^y = \mathcal{B} \backslash \mathbf{y},
    \end{equation}
    \item the mesoscale coefficients are obtained from the reduced COM-frame regression
    \begin{equation}
    \hat{\mathbf{a}}_{\mathrm{red}} = \mathbf{X}_{\mathrm{red}} \backslash \mathbf{u},
    \end{equation}
    \item and the background coefficients are obtained from the fixed-frame residual by
    \begin{equation}
    \hat{\mathbf{c}}^x = \mathcal{B} \backslash \boldsymbol{\rho}^x,
    \qquad
    \hat{\mathbf{c}}^y = \mathcal{B} \backslash \boldsymbol{\rho}^y.
    \end{equation}
\end{enumerate}
Thus the regularized formulas in
\eqref{eq:com_fast_coeff_est}, \eqref{eq:meso_regression_regularized}, and
\eqref{eq:bg_proj_closed_form} should be viewed as possible extensions of the method, not as
the path currently used by the shipped code.

The implementation is also defined on evaluated trajectory splines rather than on raw stored
sample arrays. For each drifter, the observed positions are taken as
\begin{equation}
x_{k,n} = x_k(t_{k,n}),
\qquad
y_{k,n} = y_k(t_{k,n}),
\end{equation}
and the observed velocities are taken as first derivatives of the same trajectory splines,
\begin{equation}
\dot{x}_{k,n} = \frac{d}{dt}x_k(t_{k,n}),
\qquad
\dot{y}_{k,n} = \frac{d}{dt}y_k(t_{k,n}).
\end{equation}
This keeps the fitter compatible with trajectory splines that have been smoothed before the
streamfunction fit is applied.

\subsubsection{Default spline-basis construction}

The fast temporal basis uses spline degree
\begin{equation}
S_{\mathrm{fast}} = 3
\end{equation}
by default. If no fast knot vector is supplied, the representative times
\eqref{eq:fast_knot_times} are first reduced to their unique values, and those unique
representative times are then used as the data sites for the fast spline construction.

The deployed-drifter count in the stride rule is implemented exactly as the number of drifters
whose support intervals contain the pooled time $\tau_r$, that is,
\begin{equation}
D(\tau_r)
=
\#\left\{
k : t_{k,1} \le \tau_r \le t_{k,N_k}
\right\}.
\end{equation}
The recursion \eqref{eq:fast_stride_rule} is then applied literally with that deployed count.

For the mesoscale streamfunction, the default slow basis degree is
\begin{equation}
[S_{\tilde{x}}, S_{\tilde{y}}, S_t] = [2,2,0].
\end{equation}
If no knot vectors are supplied for the mesoscale basis, the current implementation constructs
a minimal padded domain in each dimension using only repeated lower and upper endpoints. Let
\begin{equation}
q_{\min} = \min \tilde{x}_{k,n}, \qquad q_{\max} = \max \tilde{x}_{k,n},
\end{equation}
and define $r_{\min}, r_{\max}, t_{\min}, t_{\max}$ analogously. The implemented paddings are
\begin{align}
\Delta_q &= \max\!\left(50,\, 0.25(q_{\max}-q_{\min})\right),
\\
\Delta_r &= \max\!\left(50,\, 0.25(r_{\max}-r_{\min})\right),
\\
\Delta_t &= \max\!\left(1,\, 0.01(t_{\max}-t_{\min})\right).
\end{align}
The knot vectors are then formed as
\begin{align}
\Xi_q &= \left[q_{\min}-\Delta_q, \dots, q_{\min}-\Delta_q,\,
q_{\max}+\Delta_q, \dots, q_{\max}+\Delta_q\right]^\top,
\\
\Xi_r &= \left[r_{\min}-\Delta_r, \dots, r_{\min}-\Delta_r,\,
r_{\max}+\Delta_r, \dots, r_{\max}+\Delta_r\right]^\top,
\\
\Xi_t &= \left[t_{\min}-\Delta_t, \dots, t_{\min}-\Delta_t,\,
t_{\max}+\Delta_t, \dots, t_{\max}+\Delta_t\right]^\top,
\end{align}
with each lower and upper endpoint repeated $S_{\tilde{x}}+1$, $S_{\tilde{y}}+1$, and
$S_t+1$ times, respectively. This is an implementation default rather than a required part of
the mathematical model.

\subsubsection{Gauge reduction in the mesoscale solve}

The current implementation removes the spatially constant mode separately in each slow-time
block before solving for the mesoscale coefficients. If
\begin{equation}
P_{\mathrm{space}} = n_x n_y,
\qquad
L = \text{number of slow-time basis functions},
\end{equation}
then one coefficient is removed from each of the $L$ temporal blocks, leaving
\begin{equation}
L(P_{\mathrm{space}}-1)
\end{equation}
independent mesoscale unknowns.

This is done by introducing a reduced spatial lifting matrix
\begin{equation}
G_{\mathrm{space}} =
\begin{bmatrix}
I_{P_{\mathrm{space}}-1}
\\
-\mathbf{1}^\top
\end{bmatrix},
\end{equation}
so the full coefficient vector is represented as
\begin{equation}
\mathbf{a} = G \mathbf{a}_{\mathrm{red}},
\qquad
G = I_L \otimes G_{\mathrm{space}}.
\end{equation}
The regression is then solved in reduced form,
\begin{equation}
\mathbf{u} = \mathbf{X}G \mathbf{a}_{\mathrm{red}} + \boldsymbol{\eta},
\end{equation}
after which the full tensor-product coefficient vector is reconstructed by
\begin{equation}
\hat{\mathbf{a}} = G \hat{\mathbf{a}}_{\mathrm{red}}.
\end{equation}
Equivalently, the spatial coefficients within each slow-time block satisfy a zero-sum gauge.

\subsubsection{Residual-based background and retained diagnostics}

The current implementation defines the background from the fixed-frame residual
\eqref{eq:rho_fixed_residual}, not from the COM velocity \eqref{eq:com_velocity_fastcom}. In
particular,
\begin{equation}
\rho^x_{k,n} = \dot{x}_{k,n} - \hat{u}^{\mathrm{meso}}_{k,n},
\qquad
\rho^y_{k,n} = \dot{y}_{k,n} - \hat{v}^{\mathrm{meso}}_{k,n},
\end{equation}
and those residuals are projected directly onto the fast basis to obtain
$\hat{u}^{\mathrm{bg}}(t)$ and $\hat{v}^{\mathrm{bg}}(t)$. The submesoscale term is then the
remaining residual after subtracting that fitted background, exactly as in
\eqref{eq:submeso_final}.

For reproducibility of the examples and decomposition analysis, the implementation also retains
the following diagnostic quantities after the fit:
\begin{enumerate}
    \item the pooled observation times, positions, and velocities,
    \item the centered coordinates and COM velocity samples,
    \item the full and reduced mesoscale design matrices together with the reduced and lifted
    mesoscale coefficient vectors,
    \item the mesoscale, background, and submesoscale velocity components evaluated at all
    observation events.
\end{enumerate}
These diagnostics are not additional modeling assumptions, but they are part of the current
numerical realization of the method and therefore affect exact reproducibility of the MATLAB
implementation.
